\chapter{Guide de lecture et ressources ouvertes}
\label{ch:reading-guide}

Une bibliographie brute ne dit pas par où commencer. Les ressources ci-dessous
sont donc ordonnées par objectif et par niveau mathématique. Tous les intitulés
en couleur sont des liens cliquables vers une version ouverte ou vers la page
officielle de la ressource.

\section{Premiers repères statistiques}

\begin{description}[style=nextline,leftmargin=0pt]
  \item[\href{https://openintrostat.github.io/ims/}{Introduction to Modern
  Statistics --- chapitre 9}] Ressource libre de Çetinkaya-Rundel et Hardin
  \cite{openintro2024}. Le chapitre sur la régression logistique part de données,
  explique les cotes et l'interprétation des coefficients avant de développer
  l'inférence. C'est le meilleur point d'entrée de cette liste si les probabilités
  conditionnelles et les log-cotes ne sont pas encore familières.

  \item[\href{https://www.statlearning.com/}{An Introduction to Statistical
  Learning}] Ouvrage de James, Witten, Hastie, Tibshirani et Taylor
  \cite{james2023}, disponible gratuitement en versions R et Python. Le traitement
  est moins technique que \emph{The Elements of Statistical Learning}, avec des
  exemples, des figures et des laboratoires. Lire le chapitre Classification,
  puis les chapitres Resampling Methods et Linear Model Selection and
  Regularization pour comprendre validation croisée et pénalisation.
\end{description}

\section{Probabilités et calcul}

\begin{description}[style=nextline,leftmargin=0pt]
  \item[\href{https://web.stanford.edu/class/archive/cs/cs109/cs109.1264/lectures/20-LogisticRegression/}{Stanford CS109 --- Logistic Regression}]
  Cours court accompagné de diapositives et de code \cite{stanfordcs109}. Il est
  utile pour revoir maximum de vraisemblance et régression logistique dans une
  progression de probabilités appliquées, sans entrer immédiatement dans un
  traité de statistique mathématique.

  \item[\href{https://see.stanford.edu/Course/CS229}{Stanford Engineering
  Everywhere --- CS229}] Le cours historique fournit vidéos, transcriptions,
  exercices et notes. Les notes 1 couvrent classification et régression
  logistique ; les notes 5 couvrent régularisation et sélection de modèle
  \cite{ng2018cs229}. La densité mathématique est supérieure aux deux ressources
  précédentes, mais les dérivations sont directement utilisables.

  \item[\href{https://online.stat.psu.edu/stat504/}{Penn State STAT 504 ---
  Analysis of Discrete Data}] Notes de cours ouvertes du département de
  statistique de Penn State \cite{psustat504}. Les leçons 6 à 8 traitent
  successivement régression logistique binaire, diagnostics et régression
  multinomiale. Cette ressource est particulièrement utile pour approfondir
  rapports de cotes, qualité d'ajustement et interprétation statistique.
\end{description}

\section{Approfondir le formalisme}

\begin{description}[style=nextline,leftmargin=0pt]
  \item[\href{https://hastie.su.domains/ElemStatLearn/}{The Elements of
  Statistical Learning}] Référence de Hastie, Tibshirani et Friedman
  \cite{hastie2009}. Le chapitre 4 contient les équations citées dans ce document.
  Il est concis et exigeant : mieux vaut l'utiliser après une première lecture
  d'ISL ou d'OpenIntro, puis vérifier une à une vraisemblance, gradient, Hessienne
  et régularisation.

  \item[\href{https://probml.github.io/pml-book/book1.html}{Probabilistic Machine
  Learning: An Introduction}] Ouvrage de Murphy \cite{murphy2022}. Il replace la
  régression logistique dans un cadre probabiliste plus large et développe
  optimisation, calibration et modèles multiclasses. À consulter lorsque les
  objets du présent cours sont déjà maîtrisés séparément.

  \item[\href{https://www.microsoft.com/en-us/research/publication/pattern-recognition-machine-learning/}{Pattern
  Recognition and Machine Learning}] Ouvrage de Bishop \cite{bishop2006}. Sa
  section 4.3 traite les modèles linéaires de classification depuis un point de
  vue probabiliste, et développe la méthode de Newton--Raphson comme alternative
  à la descente de gradient. C'est la référence à consulter pour dépasser la
  descente du premier ordre employée ici.

  \item[\href{https://web.stanford.edu/~boyd/cvxbook/}{Convex Optimization}]
  Ouvrage ouvert de Boyd et Vandenberghe \cite{boyd2004}. Les chapitres sur la
  convexité et les méthodes de descente donnent le cadre général derrière la
  Hessienne positive et le choix d'une direction de descente. Cette référence
  traite l'optimisation mathématique générale ; son étude suppose donc déjà le
  contexte statistique de la classification.
\end{description}

\begin{resultat}[Parcours conseillé]
Pour une première étude : OpenIntro, puis ISL, puis les notes CS229. Pour vérifier
le formalisme du présent document : ESL chapitre 4 et Boyd chapitre 9. Pour aller
au-delà d'OvR vers une modélisation probabiliste multiclasses : Penn State leçon
8, puis Murphy. Cette progression évite de demander à une même source d'être à la
fois introduction, manuel de calcul et référence théorique.
\end{resultat}
